\documentclass{beamer}

% Theme choice
\usetheme{Madrid} % You can change to e.g., Warsaw, Berlin, CambridgeUS, etc.

% Encoding and font
\usepackage[utf8]{inputenc}
\usepackage[T1]{fontenc}

% Graphics and tables
\usepackage{graphicx}
\usepackage{booktabs}

% Code listings
\usepackage{listings}
\lstset{
  basicstyle=\ttfamily\small,
  keywordstyle=\color{blue},
  commentstyle=\color{gray},
  stringstyle=\color{red},
  breaklines=true,
  frame=single
}

% Math packages
\usepackage{amsmath}
\usepackage{amssymb}

% Colors
\usepackage{xcolor}

% TikZ and PGFPlots
\usepackage{tikz}
\usepackage{pgfplots}
\pgfplotsset{compat=1.18}
\usetikzlibrary{positioning}

% Hyperlinks
\usepackage{hyperref}

% Title information
\title{Chapter 2: Software Engineering \& SDLC Overview}
\author{}
\institute{}
\date{\today}

\begin{document}

\frame{\titlepage}

\begin{frame}[fragile]
    \frametitle{Introduction to Software Engineering \& SDLC}
    \begin{itemize}
        \item Software Engineering: a disciplined, systematic approach to developing, operating, and maintaining software, applying engineering principles to ensure efficiency, reliability, and maintainability.
        \item Think of it as the craftsmanship of creating complex software, similar to constructing bridges with standards.
    \end{itemize}
\end{frame}

\begin{frame}[fragile]
    \frametitle{Why is Software Engineering Important?}
    \begin{itemize}
        \item Ensures high-quality software that meets user needs.
        \item Reduces development time, cost, and risks.
        \item Facilitates teamwork among diverse professionals.
        \item Manages changes and enables scalability for the future.
    \end{itemize}
\end{frame}

\begin{frame}[fragile]
    \frametitle{Overview of the SDLC}
    \begin{itemize}
        \item The Software Development Lifecycle (SDLC) is a structured process guiding all phases of a software project.
        \item Acts as a roadmap to ensure systematic completion from concept to maintenance.
    \end{itemize}
\end{frame}

\begin{frame}[fragile]
    \frametitle{Key Phases of SDLC}
    \begin{enumerate}
        \item \textbf{Requirement Analysis}: Understand and document user needs.
        \item \textbf{Design}: Develop architecture and detailed specifications.
        \item \textbf{Development}: Implement the code based on design.
        \item \textbf{Testing}: Verify functionality, fix bugs, assure quality.
        \item \textbf{Deployment}: Launch the software for users.
        \item \textbf{Maintenance}: Support, update, and evolve the system.
    \end{enumerate}
\end{frame}

\begin{frame}[fragile]
    \frametitle{Example: Mobile Banking App}
    \begin{itemize}
        \item \textbf{Requirement Analysis}: Secure login, fund transfers, transaction history.
        \item \textbf{Design}: User interface layout, backend system architecture.
        \item \textbf{Development}: Coding features, security, functionality.
        \item \textbf{Testing}: Ensuring login security and transaction accuracy.
        \item \textbf{Deployment}: Releasing the app to customers.
        \item \textbf{Maintenance}: Security patches, new features.
    \end{itemize}
\end{frame}

\begin{frame}[fragile]
    \frametitle{Summary of Key Points}
    \begin{itemize}
        \item Software Engineering creates a disciplined framework for reliable development.
        \item SDLC provides a systematic process ensuring quality and organization.
        \item Following SDLC phases reduces errors and manages complexity.
        \item Various SDLC models exist to adapt to project needs.
    \end{itemize}
\end{frame}

\begin{frame}[fragile]
  \frametitle{Fundamental SDLC Models}
  \textbf{Overview:} \\
  The Software Development Life Cycle (SDLC) offers structured workflows guiding software projects through phases such as planning, design, implementation, testing, deployment, and maintenance. 
  \\
  Today, we focus on three core models: \textbf{Waterfall}, \textbf{Agile}, and \textbf{Spiral}, exploring their definitions, key traits, and workflows.
\end{frame}

\begin{frame}[fragile]
  \frametitle{Waterfall Model}
  \textbf{Definition:} \\
  A linear, sequential approach where each phase is completed before moving to the next, resembling a cascading waterfall.
  \\
  \textbf{Characteristics:}
  \begin{itemize}
    \item Sequential phases: Requirements $\rightarrow$ Design $\rightarrow$ Implementation $\rightarrow$ Testing $\rightarrow$ Deployment $\rightarrow$ Maintenance.
    \item Clear milestones with distinct deliverables.
    \item Emphasizes comprehensive documentation at each stage.
  \end{itemize}
  \textbf{Workflow:}
  \begin{enumerate}
    \item Gather and analyze requirements
    \item Design system architecture
    \item Develop (code) the software
    \item Test thoroughly
    \item Deploy
    \item Maintain and update
  \end{enumerate}
  \textbf{Example:} Developing a payroll system with fixed, well-understood requirements.
  \textbf{Key Point:} Best for projects with stable, well-defined specs; limited flexibility to changes.
\end{frame}

\begin{frame}[fragile]
  \frametitle{Agile Model}
  \textbf{Definition:} \\
  An iterative and incremental approach prioritizing flexibility, collaboration, and customer feedback, delivering small workable parts (increments) frequently.
  \\
  \textbf{Characteristics:}
  \begin{itemize}
    \item Iterative cycles called "sprints" (2-4 weeks)
    \item Frequent reassessment and adaptation
    \item Close stakeholder collaboration
    \item Working software as key progress indicator
  \end{itemize}
  \textbf{Workflow:}
  \begin{enumerate}
    \item Plan a sprint
    \item Develop features during the sprint
    \item Review and gather feedback
    \item Refine requirements and plan next sprint
    \item Repeat until goals achieved
  \end{enumerate}
  \textbf{Example:} Building an e-commerce site with user feedback shaping development.
  \textbf{Key Point:} Ideal for evolving or unclear requirements; requires disciplined team management.
\end{frame}

\begin{frame}[fragile]
  \frametitle{Spiral Model}
  \textbf{Definition:} \\
  Combines iterative development with risk management, emphasizing repeated cycles ("spirals") to refine the system through planning, risk analysis, engineering, and evaluation.
  \\
  \textbf{Characteristics:}
  \begin{itemize}
    \item Risk-driven: identify and mitigate risks early
    \item Iterative: each loop involves planning, analysis, design, testing
    \item Flexible to changing needs and challenges
  \end{itemize}
  \textbf{Workflow:}
  \begin{enumerate}
    \item Identify objectives, alternatives, constraints
    \item Conduct risk assessment and develop prototypes
    \item Plan next steps based on insights
    \item Develop and test system segments
    \item Repeat for refinement
  \end{enumerate}
  \textbf{Example:} Aerospace control systems needing thorough risk analysis.
  \textbf{Key Point:} Suitable for complex, high-risk projects; more costly due to iterative risk assessments.
\end{frame}

\begin{frame}[fragile]
  \frametitle{Summary Table of SDLC Models}
  \begin{tabular}{|l|l|l|l|}
    \hline
    \textbf{Model} & \textbf{Approach} & \textbf{Key Features} & \textbf{Best For} \\
    \hline
    Waterfall & Linear, Sequential & Clear milestones, strict phases & Stable, well-defined projects \\
    \hline
    Agile & Iterative, Flexible & Customer feedback, quick releases & Evolving or dynamic projects \\
    \hline
    Spiral & Risk-driven, Iterative & Risk management focus, prototyping & Complex, high-risk projects \\
    \hline
  \end{tabular}
\end{frame}

\begin{frame}[fragile]
    \frametitle{Waterfall Model}
    \begin{itemize}
        \item The Waterfall Model is an early, straightforward SDLC approach that follows a linear, sequential process.
        \item Each phase must be completed before moving to the next, flowing downhill like a waterfall.
    \end{itemize}
\end{frame}

\begin{frame}[fragile]
    \frametitle{Sequential Phases of the Waterfall Model}
    \begin{enumerate}
        \item \textbf{Requirement Analysis} -- Gather and document requirements from stakeholders.
        \item \textbf{System Design} -- Create architecture, data structures, and interface designs.
        \item \textbf{Implementation} -- Develop the software according to the design.
        \item \textbf{Integration \& Testing} -- Combine modules and verify functionality.
        \item \textbf{Deployment} -- Launch the complete system into the user environment.
        \item \textbf{Maintenance} -- Provide support, updates, and bug fixes post-deployment.
    \end{enumerate}
    \vspace{0.5cm}
    % Optional: Insert diagram as a flowchart
    % \includegraphics[width=\linewidth]{path_to_flowchart_image}
\end{frame}

\begin{frame}[fragile]
    \frametitle{Advantages of the Waterfall Model}
    \begin{itemize}
        \item \textbf{Simplicity \& Clarity} -- Easy to understand and manage due to its structured nature.
        \item \textbf{Structured Approach} -- Well-defined stages and clear deliverables.
        \item \textbf{Documentation Driven} -- Extensive documentation aids maintenance.
        \item \textbf{Early Planning} -- Requirements are defined upfront, reducing ambiguity.
    \end{itemize}
\end{frame}

\begin{frame}[fragile]
    \frametitle{Limitations of the Waterfall Model}
    \begin{itemize}
        \item \textbf{Inflexibility} -- Difficult to accommodate changes after a phase is completed.
        \item \textbf{Late Testing} -- Testing occurs late, risking costly fixes if issues are found.
        \item \textbf{Requires Clear Requirements} -- Works best when requirements are stable and well-understood.
        \item \textbf{Poor for Dynamic Projects} -- Not suitable for evolving or feedback-driven projects.
    \end{itemize}
\end{frame}

\begin{frame}[fragile]
    \frametitle{Typical Use Cases}
    Projects where requirements are well-understood and unlikely to change:
    \begin{itemize}
        \item Embedded systems
        \item Hardware development
        \item Small or straightforward applications
        \item Regulatory environments demanding thorough documentation
    \end{itemize}
    \vspace{0.5cm}
    \textbf{Key takeaway:} Suitable for projects with stable, clearly defined requirements.
\end{frame}

\begin{frame}[fragile]
    \frametitle{Summary of the Waterfall Model}
    \begin{itemize}
        \item Emphasizes a sequential, disciplined approach to SDLC.
        \item Best for projects with clear, unchanging requirements.
        \item Its rigidity may limit adaptability in fast-changing environments.
        \item Often contrasted with more iterative models like Agile.
        \item While historically significant, modern projects may prefer flexible methodologies.
    \end{itemize}
\end{frame}

\begin{frame}[fragile]
  \frametitle{Agile Model}
  \begin{itemize}
    \item An approach to software development emphasizing \textbf{iterative, incremental progress}, \textbf{flexibility}, and \textbf{stakeholder involvement}.
    \item Designed to respond quickly to changing requirements and ensure continuous delivery of valuable software.
  \end{itemize}
\end{frame}

\begin{frame}[fragile]
  \frametitle{Key Concepts of the Agile Model}
  \begin{itemize}
    \item \textbf{Iterative Development}: Break projects into small units called \textbf{sprints} (1-4 weeks), with each producing a functional piece of software. \\
    \textbf{Example}: Planning a new feature, developing it in a sprint, then demoing it to stakeholders.
    \item \textbf{Flexibility and Adaptability}: Regular reassessment of priorities allows teams to adapt swiftly to changes in requirements or market.
    \item \textbf{Stakeholder Involvement}: Continuous engagement through meetings and demos ensures alignment with stakeholder needs.
    \item \textbf{Continuous Feedback}: After each sprint, teams review progress, gather feedback, and refine for next cycles.
  \end{itemize}
\end{frame}

\begin{frame}[fragile]
  \frametitle{How Agile Works}
  \begin{enumerate}
    \item \textbf{Planning}: Prioritize features for the upcoming sprint.
    \item \textbf{Development}: Design, code, and test features within the sprint period.
    \item \textbf{Review}: Demonstrate new functionalities to stakeholders at sprint end.
    \item \textbf{Retrospective}: Reflect on the process to identify improvements.
  \end{enumerate}
\end{frame}

\begin{frame}[fragile]
    \frametitle{Spiral Model}
    \begin{itemize}
        \item A risk-driven, iterative software development process designed for complex and high-risk projects.
        \item Emphasizes repeated cycles of planning, risk analysis, engineering, and evaluation.
        \item Each iteration helps reduce uncertainties, refine features, and mitigate risks early.
    \end{itemize}
\end{frame}

\begin{frame}[fragile]
    \frametitle{Core Concepts of the Spiral Model}
    \begin{enumerate}
        \item \textbf{Risk-Driven Process}: Identify and address potential risks, such as technical feasibility or schedule delays, at each cycle.
        \item \textbf{Iterative Refinement}: Progress through multiple spiral loops, refining plans and prototypes with each iteration.
        \item \textbf{Phases of a Spiral:}
        \begin{itemize}
            \item \textbf{Planning}: Set objectives, alternatives, constraints.
            \item \textbf{Risk Analysis}: Evaluate potential problems.
            \item \textbf{Prototyping/Engineering}: Develop system parts or prototypes.
            \item \textbf{Evaluation}: Review progress and plan next steps.
        \end{itemize}
    \end{enumerate}
\end{frame}

\begin{frame}[fragile]
    \frametitle{Visualizing and Applying the Spiral}
    \begin{block}{Visual Representation}
        Imagine a series of loops expanding outward from the center, where each loop represents an iteration that adds more features, refinements, and risk mitigation.
    \end{block}
    \vspace{0.5cm}
    \textbf{Example}:
    \begin{itemize}
        \item Developing an aerospace control system with new hardware:
        \begin{itemize}
            \item \textbf{First spiral:} Prototype core algorithms and assess hardware.
            \item \textbf{Next spirals:} Refine prototypes, manage technical risks, add features.
            \item \textbf{Later spirals:} Enhance reliability, usability, regulatory compliance.
        \end{itemize}
    \end{itemize}
\end{frame}

\begin{frame}[fragile]
  \frametitle{Comparing SDLC Models}
  
  \begin{itemize}
    \item Compares Waterfall, Agile, and Spiral models across key dimensions: flexibility, risk management, user involvement, and project suitability.
  \end{itemize}
\end{frame}

\begin{frame}[fragile]
  \frametitle{Waterfall Model}
  
  \begin{itemize}
    \item \textbf{Description}: Linear, sequential process: requirements $\rightarrow$ design $\rightarrow$ implementation $\rightarrow$ testing $\rightarrow$ deployment.
    \item \textbf{Flexibility}: \textit{Low}. Changes are difficult after a phase is completed; modifications need revisiting previous stages.
    \item \textbf{Risk Management}: \textit{Limited}. Risks identified early, but late issues are costly to fix.
    \item \textbf{User Involvement}: \textit{Minimal} after initial requirements and acceptance testing.
    \item \textbf{Project Suitability}: Best for projects with well-defined, stable requirements, e.g., embedded systems with fixed specs.
  \end{itemize}
\end{frame}

\begin{frame}[fragile]
  \frametitle{Agile Model}
  
  \begin{itemize}
    \item \textbf{Description}: Iterative, incremental approach emphasizing flexibility, customer collaboration, and continuous delivery.
    \item \textbf{Flexibility}: \textit{High}. Requirements evolve through short cycles called sprints.
    \item \textbf{Risk Management}: \textit{Proactive}. Regular testing, feedback, and adaptation reduce risks early.
    \item \textbf{User Involvement}: \textit{High}. Frequent demonstrations and stakeholder interactions steer development.
    \item \textbf{Project Suitability}: Ideal for projects with changing requirements, such as web apps or startups.
  \end{itemize}
\end{frame}

\begin{frame}[fragile]
  \frametitle{Spiral Model}
  
  \begin{itemize}
    \item \textbf{Description}: Combines iterative development with risk-driven planning: each cycle involves planning, risk analysis, engineering, and evaluation.
    \item \textbf{Flexibility}: \textit{Very high}. Emphasizes incremental refinement across multiple cycles.
    \item \textbf{Risk Management}: \textit{Central}. Continuous risk assessment and mitigation at each spiral.
    \item \textbf{User Involvement}: \textit{Moderate to high}. Feedback occurs during specific cycles.
    \item \textbf{Project Suitability}: Suitable for large, complex, high-risk projects, e.g., aerospace or defense systems.
  \end{itemize}
\end{frame}

\begin{frame}[fragile]
  \frametitle{Summary and Visualization}
  
  \begin{itemize}
    \item Waterfall is simple but inflexible; best for clear requirements.
    \item Agile prioritizes adaptability and user collaboration for rapid changes.
    \item Spiral emphasizes risk management in complex, high-stakes projects.
  \end{itemize}
  
  \begin{block}{Visualization}
    \begin{center}
      \begin{tabular}{c c c}
        Waterfall & \textbf{---} & Spiral \textbf{---} & Agile \\
        (Least flexible) && (High risk management) && (Most flexible)
      \end{tabular}
    \end{center}
  \end{block}
  
  \smallskip
  \textit{Choosing the right SDLC depends on project complexity, risk, and stakeholder flexibility.}
\end{frame}

\begin{frame}[fragile]
  \frametitle{Choosing the Right SDLC Model}
  \textbf{Overview:} Selecting an appropriate Software Development Life Cycle (SDLC) model is vital for project success. The choice depends on project scope, requirements stability, team size, complexity, stakeholder involvement, and risk factors.
\end{frame}

\begin{frame}[fragile]
  \frametitle{Understanding Your Project and Requirements}
  \begin{itemize}
    \item \textbf{Well-defined & Stable Requirements} \\
    \quad - \textbf{Best Fit:} Waterfall model \\
    \quad - \textbf{Example:} Simple billing system \\
    \item \textbf{Evolving & Uncertain Requirements} \\
    \quad - \textbf{Best Fit:} Agile methodology \\
    \quad - \textbf{Example:} Mobile app with frequent updates
  \end{itemize}
\end{frame}

\begin{frame}[fragile]
  \frametitle{Team Size and Structure Considerations}
  \begin{itemize}
    \item \textbf{Small, Cross-Functional Teams:} \\
    \quad Emphasize communication; Agile works well. \\
    \item \textbf{Large, Distributed Teams:} \\
    \quad Need structured processes like Waterfall or V-Model for coordination.
  \end{itemize}
\end{frame}

\begin{frame}[fragile]
  \frametitle{Assessing Project Complexity and Risk}
  \begin{itemize}
    \item \textbf{High-Risk Projects:} \\
    \quad \textbf{Best Choice:} Spiral model \\
    \quad \textbf{Example:} Aerospace or medical device software \\
    \item \textbf{Lower Complexity Projects:} \\
    \quad Suitable for Waterfall or Agile, depending on context.
  \end{itemize}
\end{frame}

\begin{frame}[fragile]
  \frametitle{Stakeholder Involvement and Delivery Needs}
  \begin{itemize}
    \item \textbf{High User Engagement:} \\
    \quad Prefer Agile for ongoing feedback. \\
    \item \textbf{Limited User Interaction:} \\
    \quad Waterfall may be preferable when involvement is minimal past initial requirements.
  \end{itemize}
\end{frame}

\begin{frame}[fragile]
  \frametitle{Comparison of SDLC Models}
  \begin{tabular}{|l|c|c|c|}
    \hline
    \textbf{Model} & \textbf{Flexibility} & \textbf{Predictability} & \textbf{Best For} \\
    \hline
    Waterfall & Low & High & Fixed, well-understood scope \\
    \hline
    Agile & High & Less predictable & Dynamic, evolving needs \\
    \hline
    Spiral & Moderate to high & Moderate & High-risk, complex projects \\
    \hline
  \end{tabular}
\end{frame}

\begin{frame}[fragile]
  \frametitle{Additional Considerations}
  \begin{itemize}
    \item \textbf{Time to Market:} Agile delivers early and often. \\
    \item \textbf{Budget Constraints:} Fixed budgets favor Waterfall for predictability. \\
    \item \textbf{Regulations:} Strict documentation and validation favor Waterfall.
  \end{itemize}
\end{frame}

\begin{frame}[fragile]
  \frametitle{Key Takeaways}
  \begin{itemize}
    \item Match SDLC models to project characteristics for efficiency. \\
    \item Use Waterfall for stable, well-defined projects. \\
    \item Choose Agile for dynamic, rapid-changing environments. \\
    \item Spiral suits high-risk, complex projects with iterative risk management.
  \end{itemize}
\end{frame}

\begin{frame}[fragile]
  \frametitle{Application of SDLC Models}
  \textbf{Brief Summary:} Understanding which SDLC model to choose depends on project requirements, risk level, and industry context. Different models suit different scenarios, from highly structured development for safety-critical systems to flexible approaches for evolving, customer-driven projects.
\end{frame}

\begin{frame}[fragile]
  \frametitle{Waterfall Model}
  \begin{itemize}
    \item \textbf{Scenario:} Developing military or aerospace applications with fixed, well-defined requirements.
    \item \textbf{Example:} Building a spacecraft control system where specifications are set upfront and changes are costly.
    \item \textbf{Why suitable:} Ensures thorough documentation, strict regulatory compliance, and clear deliverables.
  \end{itemize}
\end{frame}

\begin{frame}[fragile]
  \frametitle{Iterative and Incremental Model}
  \begin{itemize}
    \item \textbf{Scenario:} Developing a banking system with evolving features.
    \item \textbf{Example:} Creating an online banking platform, adding modules like mobile banking in phases.
    \item \textbf{Why suitable:} Flexibility to add features iteratively, gather user feedback early, adapt easily.
  \end{itemize}
\end{frame}

\begin{frame}[fragile]
  \frametitle{V-Model}
  \begin{itemize}
    \item \textbf{Scenario:} Developing medical devices requiring rigorous validation.
    \item \textbf{Example:} Medical imaging systems where each development phase aligns with specific testing.
    \item \textbf{Why suitable:} Emphasizes verification and validation, essential for safety-critical compliance.
  \end{itemize}
\end{frame}

\begin{frame}[fragile]
  \frametitle{Spiral Model}
  \begin{itemize}
    \item \textbf{Scenario:} Large enterprise software with high risks and changing requirements.
    \item \textbf{Example:} Building a custom ERP system for a multinational, involving risk analysis and prototypes.
    \item \textbf{Why suitable:} Combines risk management with iterative development; ideal for complex, high-stakes projects.
  \end{itemize}
\end{frame}

\begin{frame}[fragile]
  \frametitle{Agile Methods}
  \begin{itemize}
    \item \textbf{Scenario:} Developing a mobile app in a startup setting.
    \item \textbf{Example:} E-commerce mobile app where customer feedback directly influences development.
    \item \textbf{Why suitable:} High flexibility, continuous involvement, quick iterations, adaptability.
  \end{itemize}
\end{frame}

\begin{frame}[fragile]
    \frametitle{Summary & Key Takeaways - Part 1}
    \textbf{Recap of Main Points on SDLC Models}
    
    \begin{itemize}
        \item \textbf{SDLC} offers a structured approach to high-quality software development.
        \item Multiple models exist, each suited for different project types and requirements.
    \end{itemize}
    
    \textbf{Major SDLC Models:}
    \begin{itemize}
        \item \textbf{Waterfall:}
        \begin{itemize}
            \item Sequential phases: Requirements $\rightarrow$ Design $\rightarrow$ Implementation $\rightarrow$ Testing $\rightarrow$ Deployment $\rightarrow$ Maintenance.
            \item Suitable for stable, well-understood requirements (e.g., government projects).
            \item Pros: Clear, manageable.
            \item Cons: Rigid, late testing.
        \end{itemize}
    \end{itemize}
\end{frame}

\begin{frame}[fragile]
    \frametitle{Summary & Key Takeaways - Part 2}
    \textbf{Major SDLC Models (cont.):}
    \begin{itemize}
        \item \textbf{V-Model:}
        \begin{itemize}
            \item Extends Waterfall with focus on testing at each stage.
            \item Ideal for safety-critical systems (medical, aerospace).
            \item Emphasizes quality assurance.
        \end{itemize}
        \item \textbf{Iterative \& Incremental:}
        \begin{itemize}
            \item Develops in cycles, adding functionality each time.
            \item Suitable for evolving requirements and early deployment (web, mobile apps).
            \item Pros: Flexibility, customer feedback.
        \end{itemize}
        \item \textbf{Spiral:}
        \begin{itemize}
            \item Combines iteration with risk assessment.
            \item Fits large, complex, high-risk projects (defense, enterprise).
            \item Focus: Risk mitigation and adaptability.
        \end{itemize}
        \item \textbf{Agile (Scrum, Kanban):}
        \begin{itemize}
            \item Focuses on adaptability, fast delivery, collaboration.
            \item Best for dynamic projects (startups, apps).
            \item Pros: Highly flexible, customer involved.
        \end{itemize}
    \end{itemize}
\end{frame}

\begin{frame}[fragile]
    \frametitle{Summary & Key Takeaways - Part 3}
    \textbf{Strategic Model Selection}
    
    \begin{itemize}
        \item Choose \textbf{Waterfall} for stable, well-defined requirements.
        \item Use \textbf{Iterative} or \textbf{Agile} when requirements may change or rapid delivery is needed.
        \item Pick \textbf{Spiral} for high-risk, complex, high-stakes projects.
        \item The \textbf{V-Model} is suited for safety-critical, verification-heavy systems.
    \end{itemize}
    
    \textbf{Example:} 
    \begin{itemize}
        \item Healthcare system with strict regulations might follow \textbf{V-Model}.
        \item A startup social media app may adopt \textbf{Agile}.
    \end{itemize}
\end{frame}

\begin{frame}[fragile]
    \frametitle{Summary & Key Takeaways - Part 4}
    \textbf{Final notes:}
    \begin{itemize}
        \item No one-size-fits-all: selection depends on project scope, complexity, and risk.
        \item Understanding strengths/weaknesses enables better process design.
        \item Strategic SDLC choice enhances resource use, quality, and stakeholder satisfaction.
    \end{itemize}
    
    \textbf{Remember:} Proper model selection is key to project success—balancing risk, flexibility, and control.
\end{frame}

\begin{frame}[fragile]
  \frametitle{Questions \& Discussion - Part 1}
  \begin{itemize}
    \item Invite questions, clarifications, and discussion to deepen understanding of SDLC concepts and model selection.
    \item Reflect on how choosing the right SDLC model impacts project success.
  \end{itemize}
\end{frame}

\begin{frame}[fragile]
  \frametitle{Questions \& Discussion - Part 2}
  \begin{block}{Key Concepts to Reinforce}
    \begin{itemize}
      \item \textbf{Why questions matter:} Help in selecting suitable SDLC models, influencing project outcomes.
      \item \textbf{Clarification areas:}
        \begin{itemize}
          \item Differences between SDLC models, e.g., Waterfall vs. Agile.
          \item Model selection criteria: project size, requirements, team, client, deadlines.
          \item Trade-offs: flexibility vs. predictability, cost, risk.
        \end{itemize}
    \end{itemize}
  \end{block}
\end{frame}

\begin{frame}[fragile]
  \frametitle{Questions \& Discussion - Part 3}
  \begin{block}{Discussion Examples}
    \begin{itemize}
      \item \textbf{Healthcare System Development:}\
      \textit{Scenario:} Developing hospital management with strict regulations. \\
      \textit{Discussion:} Is Waterfall suitable for thorough documentation or can iterative approaches adapt to evolving rules?
      
      \item \textbf{E-commerce Launch:}\
      \textit{Scenario:} Rapid rollout of a shopping platform with changing features. \\
      \textit{Discussion:} How does Agile support quick deployment and adaptability?
    \end{itemize}
  \end{block}
\end{frame}

\begin{frame}[fragile]
  \frametitle{Questions \& Discussion - Part 4}
  \begin{itemize}
    \item Emphasize aligning SDLC models with project goals. Mistakes can cause delays and increased costs.
    \item Balance flexibility (iterative) versus control (linear).
    \item Engage stakeholders early for continuous feedback, especially with Agile.
  \end{itemize}
\end{frame}

\begin{frame}[fragile]
  \frametitle{Questions \& Discussion - Part 5}
  \begin{itemize}
    \item \textbf{Discussion Prompts:} 
    \begin{itemize}
      \item "What challenges might arise with a rigid SDLC in a dynamic project?"
      \item "Share experiences where SDLC model choice affected a project’s success or failure."
    \end{itemize}
    \item Remember: The goal is to connect theory with real-world project decisions.
  \end{itemize}
\end{frame}


\end{document}